% Drop-in LaTeX corrections for the TCSSOP manuscript.
% These blocks improve the reproducibility of the computational study without
% adding the scenario graphics.

% ---------------------------------------------------------------------------
% 1) Replace the opening paragraph of Section 5.1, Experimental design
% ---------------------------------------------------------------------------

The benchmark set contains 90 scenario combinations generated from
\(10\) test matrices, \(3\) due-date scenarios, and \(3\) voltage scenarios.
Each test matrix contains 50 product projects. Hence, each benchmark instance
contains 50 projects, and the full experimental design is
\[
10 \times 3 \times 3 = 90
\]
instances. The test matrices define the binary required-test vector of each
project over the catalog in Table~\ref{tab:test_attributes_short}. The
due-date and voltage scenarios are then crossed with these test matrices to
separate the effects of portfolio composition, deadline pressure, and
voltage-compatibility intensity.

The due-date scenarios differ in both location and dispersion. Table~\ref{tab:due_scenarios}
summarizes their five-number distributions. Scenario~1 represents a relatively
concentrated due-date setting, Scenario~2 introduces a wider spread with earlier
deadlines, and Scenario~3 represents a more relaxed but still heterogeneous
deadline setting.

\begin{table}[H]
\centering
\caption{Due-date scenario summaries used in the benchmark instances}
\label{tab:due_scenarios}
\small
\begin{tabular}{lccccc}
\hline
\textbf{Due-date scenario} & \textbf{Minimum} & \textbf{Q1} & \textbf{Median} & \textbf{Q3} & \textbf{Maximum} \\
\hline
Scenario 1 & 42 & 46 & 50 & 54 & 57 \\
Scenario 2 & 30 & 40 & 49.5 & 59 & 69 \\
Scenario 3 & 34 & 43 & 53 & 63 & 69 \\
\hline
\end{tabular}
\end{table}

The voltage scenarios control the proportion of projects requiring
voltage-compatible execution. The target proportions are shown in
Table~\ref{tab:voltage_scenarios}. These settings create low, medium, and high
competition for voltage-capable chambers.

\begin{table}[H]
\centering
\caption{Voltage requirement intensity by scenario}
\label{tab:voltage_scenarios}
\small
\begin{tabular}{lcc}
\hline
\textbf{Voltage scenario} & \textbf{Description} & \textbf{Proportion of projects requiring voltage} \\
\hline
Scenario 1 & Low voltage intensity & 0.15 \\
Scenario 2 & Medium voltage intensity & 0.30 \\
Scenario 3 & High voltage intensity & 0.45 \\
\hline
\end{tabular}
\end{table}

The instances are provided to the Java implementation as a structured input
table with one row per project. Each row specifies the instance identifier,
project identifier, test-matrix scenario, due-date scenario, voltage scenario,
due date, voltage requirement indicator, and the binary required-test columns.
The implementation was developed and executed in Java using the Eclipse IDE.

% ---------------------------------------------------------------------------
% 2) Replace the implementation/parameter paragraph of Section 5.2 if desired
% ---------------------------------------------------------------------------

All algorithmic variants were implemented in Java and executed from the
Eclipse IDE under the same parameter settings. The chamber-environment
assignment model was solved by an in-house exact branch-and-bound routine over
chamber types, and the resulting chamber assignments were evaluated by the
constructive scheduler described in Section~\ref{subsec:scheduling}. No
instance-specific parameter tuning was applied.
